\section{Kの数値化}

\subsection{数値化の基本方針}

Kの定義では、Solver集合上における公理の配置を出発点として、

\[
\text{配置}
\longrightarrow
\text{配置間の関係}
\longrightarrow
\text{分離}
\]

という一続きの構造を構成した。

ここで重要なのは、Kが最終的に必要とするのは「Solver集合上に構成された公理配置の構造を、ECAの数値評価体系に接続可能な数値として表現すること」である。

したがって、数値化の段階では

\[
\boxed{
\text{公理配置の集合構造}
\longrightarrow
\text{数値}
}
\]

という変換を考える。

\subsection{数値化の対象}

公理 $a\in\mathcal{A}$ に対して、そのSolver集合上の配置は

\[
\mathcal{Z}_a = \{z\in\mathcal{Z}\mid a\in\mathcal{A}_z\}
\]

である。

また、複数公理 $A\subseteq\mathcal{A}$ に対しては、

\[
\mathcal{Z}_A = \bigcap_{a\in A}\mathcal{Z}_a
\]

であり、

\[
\mathcal{Z}_A = \{z\in\mathcal{Z}\mid A\subseteq\mathcal{A}_z\}
\]

となる。

したがって、数値化の直接の対象は、まず

\[
\mathcal{Z}_a
\]

という公理の配置集合である。

複数公理間の関係を数値化する必要がある場合には、

\[
\mathcal{Z}_a\cap\mathcal{Z}_b,
\qquad
\mathcal{Z}_a\setminus\mathcal{Z}_b,
\qquad
\mathcal{Z}_b\setminus\mathcal{Z}_a,
\qquad
\mathcal{Z}_a\triangle\mathcal{Z}_b
\]

などの集合構造を対象とすることができる。

ただし、Kの基本的な数値化において、これらすべてを一つの式へ同時に組み込む必要はない。

Kの中心は、まず各公理のSolver集合上における配置である。

\subsection{最も単純な数値化}

Solver集合 $\mathcal{Z}$ が有限集合である場合、最も直接的な数値化として、

\[
\boxed{
K(a) = \frac{|\mathcal{Z}_a|}{|\mathcal{Z}|}
}
\]

を考えることができる。

ここで、

\[
\mathcal{Z}_a = \{z\in\mathcal{Z}\mid a\in\mathcal{A}_z\}
\]

である。

この値は、

\[
0\leq K(a)\leq1
\]

を満たす。

また、

\[
K(a)=0
\]

ならば、

\[
\mathcal{Z}_a=\varnothing
\]

であり、公理 $a$ は対象とするSolver集合上には配置されていない。

一方、

\[
K(a)=1
\]

ならば、

\[
\mathcal{Z}_a=\mathcal{Z}
\]

であり、公理 $a$ は対象となるすべてのSolverに含まれている。

したがって、

\[
K(a) = \frac{|\mathcal{Z}_a|}{|\mathcal{Z}|}
\]

は、公理 $a$ がSolver集合全体のどの程度に配置されているかを表す最も単純な数値として解釈できる。

\subsection{具体例}

例えば、

\[
\mathcal{Z} = \{z_1,z_2,z_3,z_4\}
\]

とする。

公理 $a$ が

\[
\mathcal{Z}_a = \{z_1,z_2\}
\]

に配置されているなら、

\[
|\mathcal{Z}_a|=2,
\qquad
|\mathcal{Z}|=4
\]

であるため、

\[
K(a) = \frac{2}{4} = 0.5
\]

となる。

同様に、

\[
\mathcal{Z}_b = \{z_3,z_4\}
\]

ならば、

\[
K(b) = \frac{2}{4} = 0.5
\]

である。

しかし、

\[
\mathcal{Z}_a\cap\mathcal{Z}_b = \varnothing
\]

であるため、両者はSolver集合上では完全に分離している。

この例から、

\[
K(a)=K(b)
\]

であっても、

\[
\mathcal{Z}_a=\mathcal{Z}_b
\]

とは限らないことが分かる。

したがって、

\[
\frac{|\mathcal{Z}_a|}{|\mathcal{Z}|}
\]

は配置の「量」を数値化するものであり、配置そのものの完全な構造を一つの実数へ圧縮するものではない。

これは欠点ではなく、Kと集合構造の役割を区別するために重要な性質である。

\subsection{配置構造と数値の区別}

Kの数値化においては、

\[
\mathcal{Z}_a
\]

と

\[
K(a)
\]

を同一視してはならない。

$\mathcal{Z}_a$ は公理 $a$ のSolver集合上における配置そのものであり、

\[
K(a)
\]

はその配置を数値として評価した結果である。

したがって、

\[
\boxed{
\mathcal{Z}_a
\neq
K(a)
}
\]

である。

より正確には、

\[
\mathcal{Z}_a
\longrightarrow
K(a)
\]

という数値評価が行われる。

また、すべての公理について考えれば、

\[
\{\mathcal{Z}_a\mid a\in\mathcal{A}\}
\]

がSolver集合上の公理配置全体を与え、それに対する数値評価として

\[
\{K(a)\mid a\in\mathcal{A}\}
\]

が得られる。

したがって、

\[
\boxed{
\{\mathcal{Z}_a\mid a\in\mathcal{A}\}
\longrightarrow
{K(a)\mid a\in\mathcal{A}}
}
\]

という対応を考えることができる。

\subsection{配置間の関係を含む場合}

Kの定義では、公理間の関係および分離は各公理の配置から導かれる。

例えば、

\[
\mathcal{Z}_a\cap\mathcal{Z}_b
\]

は共通配置を表し、

\[
\mathcal{Z}_a\triangle\mathcal{Z}_b
\]

は配置の差を表す。

したがって、必要に応じてこれらの集合に対しても数値評価を行うことができる。

例えば有限なSolver集合において、

\[
\frac{|\mathcal{Z}_a\cap\mathcal{Z}_b|}{|\mathcal{Z}|}
\]

は、両公理が共通して配置されているSolverの割合を表す。

また、

\[
\frac{|\mathcal{Z}_a\triangle\mathcal{Z}_b|}{|\mathcal{Z}|}
\]

は、両公理の配置が異なるSolverの割合を表す。

完全分離の場合には、

\[
\mathcal{Z}_a\cap\mathcal{Z}_b=\varnothing
\]

であるため、

\[
|\mathcal{Z}_a\cap\mathcal{Z}_b|=0
\]

となる。

ただし、これらをKの基本式へ必ず組み込む必要はない。

これらは、

\[
\mathcal{Z}_a
\longrightarrow
\mathcal{Z}_a\text{ と }\mathcal{Z}_b\text{ の関係}
\]

という構造を数値的に調べるための補助的な評価として位置付けることができる。

Kの中心はあくまで公理のSolver集合上の配置であり、二つの公理間の距離をKそのものとして定義することではない。

\subsection{ECAの既存の数値評価との接続}

ECAでは、公理体系に対する数値評価について、まず非正規化された値を定め、それを必要に応じて正規化する構造が既に導入されている。

一般に、

\[
\widetilde\mu(i) = W(i)
\]

として非正規化された値を定め、有限かつ正の上限が存在する場合には、

\[
\mu_{\mathrm{global}}(i) = \frac{\widetilde\mu(i)}{\sup_{j\in\mathcal{F}_0}\widetilde\mu(j)}
\]

によって、

\[
\mu_{\mathrm{global}}(i)\in[0,1]
\]

へ正規化できる。

Kについても、同じ考え方を利用できる。

すなわち、Solver集合上の公理配置 $\mathcal{Z}_a$ に対して、まず非正規化された数値

\[
\widetilde{K}(a)
\]

を与え、その後必要に応じて正規化する。

一般形として、

\[
\widetilde{K}(a) = W_K(\mathcal{Z}_a)
\]

と置けば、

\[
K(a) = \frac{\widetilde{K}(a)}{\sup_{b\in\mathcal{F}_K}\widetilde{K}(b)}
\]

という正規化を考えることができる。

ここで $\mathcal{F}_K$ はKを評価する対象となる公理の集合である。

上限が有限かつ正であれば、

\[
K(a)\in[0,1]
\]

となる。

この形式は、Kのために新たな数値評価原理を導入するのではなく、ECAですでに使用している「非正規化された値を構成し、必要に応じて正規化する」という数値評価の考え方を、Solver集合上の公理配置へ適用するものである。

\subsection{有限Solver集合における単純化}

有限なSolver集合

\[
\mathcal{Z}
\]

に対して、公理の各Solverへの寄与を一様に扱う場合、

\[
W_K(\mathcal{Z}_a) = |\mathcal{Z}_a|
\]

と置くことができる。

この場合、

\[
\widetilde{K}(a) = |\mathcal{Z}_a|
\]

となる。

さらに、

\[
\sup_{b\in\mathcal{F}_K}|\mathcal{Z}_b|
\]

を基準として正規化することもできる。

一方、Solver集合そのものを基準とする割合として評価する場合には、

\[
\boxed{K(a) = \frac{|\mathcal{Z}_a|}{|\mathcal{Z}|}}
\]

となる。

したがって、

\[
K(a) = \frac{|\mathcal{Z}_a|}{|\mathcal{Z}|}
\]

は、Kの定義から独立した新しい原理として突然導入されるのではなく、有限Solver集合に対して一様な数え上げを採用した場合の具体的な数値化として位置付けられる。

\subsection{数値化における「分散」の扱い}

Kにおける「分散」は

\[
\text{公理がSolver集合上のどこに配置されているか}
\]

という構造を基礎としている。

例えば

\[
\mathcal{Z}_a = \{z_1,z_2\}
\]

と

\[
\mathcal{Z}_b = \{z_3,z_4\}
\]

がある場合、

\[
|\mathcal{Z}_a| = |\mathcal{Z}_b|
\]

であっても、

\[
\mathcal{Z}_a\cap\mathcal{Z}_b = \varnothing
\]

である。

したがって、単純な数値

\[
K(a)=K(b)
\]

だけでは配置関係まで復元できない。

このことからも、Kの「分散」はまず集合論的な配置構造を構成し、その構造を数値として評価するものとして扱うことが分かる。

\subsection{無限Solver集合への拡張}

$\mathcal{Z}$ が無限集合の場合には、

\[
\frac{|\mathcal{Z}_a|}{|\mathcal{Z}|}
\]

を通常の有限集合と同じ意味で実数値として扱うことはできない場合がある。

特に、

\[
|\mathcal{Z}_a| = \aleph_0
\]

で

\[
|\mathcal{Z}| = \aleph_0
\]

のような場合、

\[
\frac{\aleph_0}{\aleph_0}
\]

は、有限集合における割合のような実数値を与えない。

したがって、無限Solver集合については、ECAですでに導入されている数値評価・正規化の枠組みを利用し、対象となるSolver集合に対して適切な非正規化値を構成する必要がある。

一般形としては、

\[
\widetilde{K}(a) = W_K(\mathcal{Z}_a)
\]

とし、

\[
K(a) = \frac{\widetilde{K}(a)}{\sup_{b\in\mathcal{F}_K}\widetilde{K}(b)}
\]

という正規化を用いることができる。

さらに、対象となるSolver集合に自然な測度が導入されている場合には、

\[
\widetilde{K}(a)=\mu(\mathcal{Z}_a)
\]

として、

\[
K(a) = \frac{\mu(\mathcal{Z}_a)}{\mu(\mathcal{Z})}
\]

のような評価も可能になる。

ただし、これは「Kは必ず測度によって定義されなければならない」という意味ではない。

Kに必要なのは最終的に数値であり、測度はその数値化を実現する一つの方法である。

したがって、

\[
\boxed{
\text{Kの本質は測度ではなく、Solver集合上の公理配置を数値化すること}
}
\]

である。

\subsection{Kの数値と配置構造の役割分担}

Kの数値化では、$\mathcal{Z}_a$ という「集合構造」と $K(a)$ という「数値」を明確に区別する。

集合構造は配置そのものを保持するが、一方で数値は「その配置についてECAの数値評価へ渡すための情報」を与える。

したがって、

\[
\mathcal{Z}_a
\longrightarrow
\widetilde{K}(a)
\longrightarrow
K(a)
\]

という段階を考えることができる。

有限かつ一様な場合には、

\[
\mathcal{Z}_a
\longrightarrow
|\mathcal{Z}_a|
\longrightarrow
\frac{|\mathcal{Z}_a|}{|\mathcal{Z}|}
\]

となる。

この構造により、$\mathcal{Z}_a$ が配置の構造を担い、$K(a)$ がその数値評価を担うという役割分担が明確になる。

\subsection{B,I,Kによる総合評価}

Kは公理配置に関する数値であるが、ECAにおける最終的な数値評価では、K単独で全体を表現する必要はない。

ECAでは、

\[
T_{\mathrm{overlay}} = \frac{1}{N}\sum_{i=1}^{N}B_iI_iK_i
\]

によって、B、I、Kを組み合わせた評価を行う。

ここで、B、I、Kはそれぞれ次の要素を担う。

\begin{enumerate}
  \item $B_i$ : 公理数や公理選択空間に関する側面
  \item $I_i$ : 構造整合性
  \item $K_i$ : Solver集合上の公理配置
\end{enumerate}

したがって、$K_i$ だけでは表現しきれないものについても、$B_i$, $I_i$, $K_i$ という組み合わせによって評価することができる。
このため、Kの数値化に対して「K単独の数値から公理配置の全構造を復元できなければならない」という条件を要求する必要はない。

\subsection{数値の精度}

Kを実際に数値として算出する場合には、内部計算と表示上の丸めを区別する。

例えば、

\[
K(a)=\frac{|\mathcal{Z}_a|}{|\mathcal{Z}|}
\]

から得られた値が有限小数でない場合でも、計算段階では可能な限り元の値を保持し、表示段階で必要な桁数へ丸めることができる。

例えば、

\[
K(a)=\frac{17}{37}
\]

であれば、

\[
K(a) = 0.459459459459\ldots
\]

となる。

内部計算では、

\[
\frac{17}{37}
\]

という値を保持し、表示時には必要な精度に応じて

\[
0.4594594595
\]

などと表現することができる。
必要に応じて切り上げや丸めを行う場合、それは数値評価後の表示規則として扱い、Kの基礎となる数値そのものとは区別する。

\subsection{数値化の全体構造}

以上をまとめると、Kの数値化は次の順序で整理できる。

まず、各公理 $a$ に対して、

\[
\mathcal{Z}_a = \{z\in\mathcal{Z}\mid a\in\mathcal{A}_z\}
\]

を構成する。

次に、必要に応じて、

\[
\mathcal{Z}_a\cap\mathcal{Z}_b,
\]

\[
\mathcal{Z}_a\setminus\mathcal{Z}_b,
\]

\[
\mathcal{Z}_b\setminus\mathcal{Z}_a,
\]

\[
\mathcal{Z}_a\triangle\mathcal{Z}_b
\]

などによって配置間の関係を調べる。

その上で、配置を数値化するための非正規化値

\[
\widetilde{K}(a) = W_K(\mathcal{Z}_a)
\]

を構成する。

必要に応じて、

\[
K(a) = \frac{\widetilde{K}(a)}{\sup_{b\in\mathcal{F}_K}\widetilde{K}(b)}
\]

によって正規化する。

有限Solver集合を一様に数え上げる最も単純な場合には、

\[
\boxed{K(a) = \frac{|\mathcal{Z}_a|}{|\mathcal{Z}|}}
\]

となる。

したがって、Kの数値化全体は、

\[
\boxed{
\mathcal{Z}_a
\longrightarrow
\text{配置構造}
\longrightarrow
\widetilde{K}(a)
\longrightarrow
K(a)
}
\]

として表現できる。

さらにECA全体の数値評価へ接続する場合には、

\[
\boxed{T_{\mathrm{overlay}} = \frac{1}{N}\sum_{i=1}^{N}B_iI_iK_i}
\]

となる。

このとき、Kは

\[
\boxed{
\text{Solver集合上における公理配置の数値評価}
}
\]

として位置付けられる。

\subsection{Kの数値化における確定事項}

以上のことから、Kの数値化について現時点で確定している構造は次の通りである。

\begin{enumerate}
\item Kの基礎対象は、各Solverに付随する公理構造 $\mathcal{A}_z$ であること。
\item 公理 $a$ のSolver集合上の配置は $\mathcal{Z}_a = \{z\in\mathcal{Z}\mid a\in\mathcal{A}_z\}$ によって表現すること。
\item 複数公理の共通配置は $\mathcal{Z}_A = \bigcap_{a\in A}\mathcal{Z}_a$ によって表現すること。
\item 配置間の関係は、共通部分、差集合、包含関係、対称差などの集合演算によって導出すること。
\item 分離は、$\mathcal{Z}_a\cap\mathcal{Z}_b=\varnothing$ のように、配置間の関係から導かれること。
\item Kはこれらの集合構造そのものではなく、それらを数値化した量であること。
\item 有限Solver集合を一様に扱う最も単純な具体化として、$K(a) = \frac{|\mathcal{Z}_a|}{|\mathcal{Z}|}$ を用いることができること。
\item 無限Solver集合など、単純な要素数の比が使えない場合には、ECA既存の非正規化値と正規化の枠組みによって数値化する。
\item KはSolver集合上の公理配置を数値として評価する量であること。
\item K単独ですべての公理体系情報を表現する必要はなく、BおよびIと組み合わせて $T_{\mathrm{overlay}} = \frac{1}{N} \sum_{i=1}^{N}B_iI_iK_i$ によって総合的に評価すること。
\end{enumerate}

以上により、Kは集合論的なSolver構造から出発し、その配置構造を経てECAの既存の数値評価体系へ接続される量として定式化できる。